\begin{tabularx}{\textwidth}{L{31mm}L{44mm}X}
\toprule
\textbf{Fogalom} & \textbf{Definíció / jelölés} & \textbf{Értelmezés} \\
\midrule
Időigény & $T_M(n)$ a legrosszabb futási idő $n$ hosszú bemeneteken. & A bemenethossz függvényében mérjük, nem egy konkrét bemenetre. \\
Polinomiális idő & $T_M(n)=O(n^c)$ valamely konstans $c$ mellett. & Elméleti értelemben hatékonynak tekintett algoritmus. \\
$\mathrm{P}$ osztály & Polinomiális időben eldönthető nyelvek osztálya. & $\mathrm{P}\subseteq\mathrm{R}$, mert polinomiális időben futó eldöntő gép biztosan megáll. \\
Exponenciális idő & Például $O(2^n)$ vagy $O(c^n)$. & Sok kombinatorikus keresési probléma naiv megoldása ilyen. \\
Kódolás szerepe & A bemenet hossza a leírás hosszát jelenti. & Egy szám értéke és a számjegyeinek száma nagyon eltérő lehet, ezért fontos a bitméret szerinti elemzés. \\
\bottomrule
\end{tabularx}
